\section{Conclusiones}

El núcleo de Eriksen proporciona una base útil para representar el
\co{} larvario, pero su transferencia directa no describe adecuadamente
los ocho cierres propios. Las discrepancias son compatibles con la
ausencia de una fuente microbiana explícita, la transferencia de
parámetros a otro dominio experimental y la falta de acople entre
alimento y asimilación.

En calibración completa, el modelo acoplado redujo el RMSE censurado
conjunto de $183{,}61$ a $85{,}16$ ppm/min, una disminución de
$53{,}6$ %. En validación cruzada interna, la reducción puntual fue de
$31{,}95$ %, con un RMSE de $124{,}94$ ppm/min. No obstante, el modelo
acoplado no mejoró el RMSE conjunto de los cierres exactos ni la
cobertura de las cotas saturadas. El bootstrap exploratorio incluyó
reducciones negativas, por lo que no se demuestra una superioridad
predictiva concluyente.

La descomposición controlada atribuyó contablemente $55{,}8$ % de la
mejora de calibración a la recalibración del núcleo, $12{,}7$ % a la
fuente microbiana aditiva y $31{,}5$ % al acople con el alimento. Estos
valores no representan efectos causales independientes. El acumulado de
\co{} fue relativamente estable ante perturbaciones locales y dentro de
las regiones de costo casi equivalente; sin embargo, $a_{\max}$ alcanzó
la frontera superior y el límite superior de $\rho$ no quedó bien
identificado.

La evaluación externa muestra mejoras específicas, pero también fallos
en la calidad dietaria, la ubicación de picos, la partición microbiana y
la respuesta a temperaturas extremas. En consecuencia, el modelo se
presenta como una extensión local de Eriksen con mayor capacidad de
ajuste y una mejoría predictiva interna puntual, no como un predictor
de alta certeza ni como un sustituto universalmente superior.
